% Кодировка и язык
\usepackage{fontspec}
\usepackage[english,russian]{babel}

\setmainfont{Times New Roman}
\setmonofont[
    Scale=MatchLowercase
]{Noto Sans Mono}

% Поля страницы
\usepackage[
    left=30mm,
    right=15mm,
    top=20mm,
    bottom=20mm
]{geometry}

% Межстрочный интервал и абзацы
\usepackage{setspace}
\onehalfspacing
\setlength{\parindent}{1.25cm}

\usepackage{enumitem}

\setlist[itemize]{
    itemsep=0pt,
    topsep=0.3em,
    parsep=0pt,
    partopsep=0pt
}

% Математика
\usepackage{amsmath,amssymb}
\newtheorem{definition}{Определение}[section]
\newtheorem{theorem}{Теорема}[section]
\newtheorem{lemma}{Лемма}[section]

% Рисунки
\usepackage{graphicx}
\usepackage{float}

% Заголовки
\usepackage{titlesec}
\usepackage{tocloft}
\setcounter{secnumdepth}{3}
\setcounter{tocdepth}{3}

\renewcommand{\cftsecaftersnum}{.}
\renewcommand{\cftsubsecaftersnum}{.}

\setlength{\cftsecnumwidth}{2em}
\setlength{\cftsubsecnumwidth}{3em}
\renewcommand{\thesection}{\arabic{section}}
\renewcommand{\thesubsection}{\thesection.\arabic{subsection}}

\titleformat{\section}
{\bfseries\Large\centering}
{\thesection.}
{0.5em}
{}

\titleformat{\subsection}
{\bfseries\large}
{\thesubsection.}
{0.5em}
{}

% Таблицы
\usepackage{array}
\usepackage{booktabs}
\usepackage{multirow}
\usepackage{tabularx}
\usepackage{tabularray}

% Подписи
\usepackage{caption}
\usepackage{subcaption}

\captionsetup{
    justification=centering,
    labelsep=period
}

\captionsetup[table]{
    font=small
}

% Цвета
\usepackage{xcolor}

% Код Python
% --- 5. ОФОРМЛЕНИЕ КОДА (MINTED) ---
\usepackage[outputdir=build]{minted}

\usepackage{tcolorbox}
\definecolor{bg}{gray}{0.95}
\tcbuselibrary{minted,breakable,xparse,skins}
\DeclareTCBListing{mintedbox}{O{}m!O{}}{%
  breakable=true,
  listing engine=minted,
  listing only,
  minted language=#2,
  minted style=default,
  minted options={linenos,gobble=0,breaklines=true,breakafter=,,fontsize=\small,numbersep=8pt,#1},
  boxsep=0pt,
  left=25pt, right=0pt, top=3pt, bottom=3pt,
  arc=5pt,
  colback=bg,
  colframe=orange!70,
  enhanced,
  overlay={%
    \begin{tcbclipinterior}
    \fill[orange!20!white] (frame.south west) rectangle ([xshift=20pt]frame.north west);
    \end{tcbclipinterior}},
  #3}

% Блок с кодом из внешнего файла в оформлении mintedbox.
\newcommand{\inputmintedbox}[3][]{%
    \begin{tcolorbox}[
        breakable,
        enhanced,
        boxsep=0pt,
        left=25pt,
        right=0pt,
        top=3pt,
        bottom=3pt,
        arc=5pt,
        colback=bg,
        colframe=orange!70,
        overlay={%
            \begin{tcbclipinterior}
            \fill[orange!20!white]
                (frame.south west)
                rectangle ([xshift=20pt]frame.north west);
            \end{tcbclipinterior}
        }
    ]
        \inputminted[
            linenos,
            breaklines,
            fontsize=\small,
            numbersep=8pt,
            #1
        ]{#2}{#3}
    \end{tcolorbox}%
}

% Оформление консольных сессий. Лексер linuxconsole и стиль linuxterm
% поставляются локальным пакетом terminal_theme (см. README.md).
\definecolor{termbg}{HTML}{1E1E1E}
\definecolor{termfg}{HTML}{D4D4D4}
\definecolor{termbar}{HTML}{2D2D2D}
\definecolor{termframe}{HTML}{3C3C3C}
\definecolor{termaccent}{HTML}{4EC9B0}

\DeclareTCBListing{consolebox}{O{}}{
    enhanced,
    breakable,
    listing only,
    listing engine=minted,
    minted language=linuxconsole,
    minted style=linuxterm,
    minted options={
        linenos=false,
        breaklines=true,
        fontsize=\small
    },
    colback=termbg,
    coltext=termfg,
    colframe=termframe,
    boxrule=0.5pt,
    arc=3mm,
    boxsep=0pt,
    left=4mm,
    right=4mm,
    top=3mm,
    bottom=3mm,
    colbacktitle=termbar,
    coltitle=termaccent,
    fonttitle=\ttfamily\small\bfseries,
    title={
        \textcolor{red!75}{\textbullet}\,
        \textcolor{yellow!80}{\textbullet}\,
        \textcolor{green!65}{\textbullet}
        \hspace{2mm} TERMINAL
    },
    titlerule=0pt,
    toptitle=2mm,
    bottomtitle=2mm,
    before skip=10pt,
    after skip=10pt,
    #1
}

\usepackage{relsize}

% Ссылки
\usepackage[hidelinks]{hyperref}
